%%%%%%%%%%%%%%%%%%%%%%%%%%%%%%%%%%%%%%%%%%%%%%%%%%%%%%%%%%%%%%%%%%%%%%%%%%%
\chapter{Standard Relativistic Quantum Mechanics and Its Limitations}

\epigraph{
``The difficulty is not to construct relativistic wave equations.
The difficulty is to understand what they mean.''
}{Author's paraphrase of the historical development}

%%%%%%%%%%%%%%%%%%%%%%%%%%%%%%%%%%%%%%%%%%%%%%%%%%%%%%%%%%%%%%%%%%%%%%%%%%%
\section{Historical Introduction and the Axioms of Relativistic Quantum Mechanics}

\subsection{3.1 Historical Perspective}

Relativistic quantum mechanics (RQM) arose from an attempt to reconcile the
two most successful physical theories of the early twentieth century:
quantum mechanics and special relativity.  The development was remarkably
rapid. Einstein's formulation of special relativity (1905) established the
Lorentz invariance of the laws of physics, while the matrix and wave
formulations of quantum mechanics (1925--1926) provided a consistent
description of microscopic phenomena. It was therefore natural to seek a
single-particle wave equation respecting both principles.

The first successful relativistic wave equation was proposed independently by
Klein, Gordon and Fock during 1926--1927,

\begin{equation}
\left(
\Box+\frac{m^2c^2}{\hbar^2}
\right)\psi=0,
\label{eq:KG}
\end{equation}

where

\begin{equation}
\Box
=
\frac1{c^2}\frac{\partial^2}{\partial t^2}
-\nabla^2
\end{equation}

is the d'Alembert operator.

Equation (\ref{eq:KG}) is Lorentz invariant and correctly reproduces the
relativistic energy-momentum relation,

\begin{equation}
E^2
=
p^2c^2+m^2c^4.
\label{eq:EinsteinEnergy}
\end{equation}

Nevertheless, it immediately encountered serious conceptual difficulties.
Unlike the Schrödinger equation, it is second order in time and possesses
solutions of both positive and negative frequency,

\begin{equation}
E=\pm\sqrt{p^2c^2+m^2c^4}.
\end{equation}

More importantly, the conserved current derived from
(\ref{eq:KG}) cannot generally be interpreted as a probability current,
because its temporal component is not positive definite.

These problems motivated Dirac's search for a first-order relativistic wave
equation. His celebrated equation,

\begin{equation}
(i\hbar\gamma^\mu\partial_\mu-mc)\psi=0,
\label{eq:Dirac}
\end{equation}

introduced spin naturally, predicted the existence of antiparticles,
and successfully explained the fine structure of hydrogen.

Despite these achievements, even the Dirac equation failed to provide a fully
satisfactory relativistic one-particle theory. The difficulties were no longer
merely technical but conceptual.

Among them are

\begin{enumerate}

\item absence of a covariant position operator,

\item impossibility of strict localization,

\item negative-energy solutions,

\item pair creation,

\item incompatibility of fixed particle number with relativistic interactions,

\item difficulties in constructing relativistic many-body systems.

\end{enumerate}

These problems eventually led to quantum field theory, in which fields rather
than particles become the fundamental dynamical variables.

The purpose of the present chapter is not to present quantum field theory.
Instead, we ask a different question.

\begin{quote}
How far can one proceed while insisting that the fundamental object is still a
single relativistic quantum particle?
\end{quote}

Answering this question is essential for understanding why alternative
approaches---algebraic quantum field theory, modular localization,
string-localized fields, direct interaction theories, and
Stueckelberg--Horwitz--Piron mechanics---adopt different notions of
localization and particle ontology.

%%%%%%%%%%%%%%%%%%%%%%%%%%%%%%%%%%%%%%%%%%%%%%%%%%%%%%%%%%%%%

\subsection{3.2 The Logical Structure of Relativistic Quantum Mechanics}

Ordinary quantum mechanics rests upon several mathematical assumptions that
are largely independent of the specific Hamiltonian.  Relativistic quantum
mechanics attempts to preserve these assumptions while replacing Galilean
symmetry by Poincaré symmetry.

The theory may therefore be formulated through the following axioms.

\paragraph{Axiom I (Hilbert Space).}

The physical states form a complex separable Hilbert space

\begin{equation}
\mathcal H,
\end{equation}

equipped with a positive-definite inner product

\begin{equation}
(\psi,\phi).
\end{equation}

Physical states correspond to rays in $\mathcal H$.

\paragraph{Axiom II (Observables).}

Every measurable quantity is represented by a self-adjoint operator acting on
$\mathcal H$.

Expectation values are

\begin{equation}
\langle A\rangle
=
(\psi,A\psi).
\end{equation}

\paragraph{Axiom III (Probability).}

The inner product determines probabilities according to Born's rule.

For every projection operator $P$,

\begin{equation}
\mathrm{Prob}(P)
=
(\psi,P\psi).
\end{equation}

This axiom immediately implies that probabilities must remain non-negative.

%%%%%%%%%%%%%%%%%%%%%%%%%%%%%%%%%%%%%%%%%%%%%%%%%%%%%%%%%%%%%%%

\paragraph{Axiom IV (Relativistic Symmetry).}

The Hilbert space carries a unitary representation of the universal covering
group of the proper orthochronous Poincaré group,

\begin{equation}
U(a,\Lambda).
\end{equation}

Its infinitesimal generators satisfy

\begin{align}
[P^\mu,P^\nu]
&=0,
\\
[M^{\mu\nu},P^\rho]
&=
i\hbar
(g^{\nu\rho}P^\mu-g^{\mu\rho}P^\nu),
\\
[M^{\mu\nu},M^{\rho\sigma}]
&=
i\hbar
(
g^{\mu\sigma}M^{\nu\rho}
+
g^{\nu\rho}M^{\mu\sigma}
-
g^{\mu\rho}M^{\nu\sigma}
-
g^{\nu\sigma}M^{\mu\rho}
).
\end{align}

The Casimir operators are

\begin{align}
P^\mu P_\mu &=m^2c^2,
\\
W^\mu W_\mu
&=
-m^2c^2s(s+1),
\end{align}

where

\begin{equation}
W^\mu
=
\frac12
\epsilon^{\mu\nu\rho\sigma}
P_\nu
M_{\rho\sigma}
\end{equation}

is the Pauli--Lubanski vector.

Thus mass and spin label the irreducible particle representations.

%%%%%%%%%%%%%%%%%%%%%%%%%%%%%%%%%%%%%%%%%%%%%%%%%%%%%%%%%%%%%%%

\paragraph{Axiom V (Positive Energy).}

The Hamiltonian satisfies

\begin{equation}
P^0\ge0.
\end{equation}

The vacuum is therefore stable.

Although seemingly innocuous, this requirement ultimately lies behind
Hegerfeldt's theorem, the Reeh--Schlieder theorem, and several no-go results
discussed in the previous chapter.

%%%%%%%%%%%%%%%%%%%%%%%%%%%%%%%%%%%%%%%%%%%%%%%%%%%%%%%%%%%%%%%

\paragraph{Axiom VI (Localization).}

One wishes to associate with every spatial region $\Omega$ a projection
operator

\begin{equation}
E(\Omega),
\end{equation}

interpreted as localization inside $\Omega$.

The existence of such projectors is automatic in nonrelativistic quantum
mechanics.

Whether they exist in relativistic quantum mechanics is one of the central
questions of this monograph.

%%%%%%%%%%%%%%%%%%%%%%%%%%%%%%%%%%%%%%%%%%%%%%%%%%%%%%%%%%%%%%%

\subsection{3.3 Why These Axioms Are Difficult to Reconcile}

Each axiom appears natural when considered separately.

Together, however, they generate deep tensions.

Positive energy requires the Hamiltonian

\begin{equation}
H
=
\sqrt{m^2c^4+c^2P^2},
\end{equation}

which is a nonlocal pseudodifferential operator.

Lorentz covariance mixes space and time in a manner incompatible with the
spectral theory of an ordinary position operator.

Strict localization conflicts with analyticity properties implied by the
positive-energy spectrum.

Finally, relativistic interactions inevitably allow particle creation and
annihilation, contradicting the assumption that the Hilbert space describes a
fixed number of particles.

Consequently, the remainder of this chapter is devoted not merely to deriving
relativistic wave equations but to identifying precisely where these axioms
cease to be mutually compatible.

%%%%%%%%%%%%%%%%%%%%%%%%%%%%%%%%%%%%%%%%%%%%%%%%%%%%%%%%%%%%%%%

\subsection{3.4 Roadmap of the Chapter}

The discussion proceeds in the following logical order.

\begin{enumerate}

\item The Klein--Gordon equation is derived from the relativistic dispersion
relation.

\item Its conserved current and probability interpretation are analyzed.

\item Dirac's linearization procedure is developed from first principles.

\item Spinors and Clifford algebras are introduced.

\item The Foldy--Wouthuysen representation is constructed.

\item Localization of relativistic wave packets is examined.

\item Pair creation and the Klein paradox are analyzed.

\item Relativistic many-body mechanics is considered.

\item The limitations of fixed-particle-number quantum mechanics are
summarized.

\end{enumerate}

Only after these limitations have been established shall we proceed to quantum
field theory and to alternative relativistic frameworks that abandon the
assumption of point-localized particles.
%%%%%%%%%%%%%%%%%%%%%%%%%%%%%%%%%%%%%%%%%%%%%%%%%%%%%%%%%%%%%%%%%%%%%%%%%%
\chapter{Quantum Field Theory and Localization}

\epigraph{
``A quantum field is not a collection of particles.
Particles are certain states of a quantum field.''
}{Adapted from the modern viewpoint}

%%%%%%%%%%%%%%%%%%%%%%%%%%%%%%%%%%%%%%%%%%%%%%%%%%%%%%%%%%%%%%%%%%%%%%%%%%

\section{Introduction}

The transition from relativistic quantum mechanics to quantum field theory is
often presented as a consequence of particle creation and annihilation.
Although this explanation is historically accurate, it is conceptually
incomplete.

The more fundamental reason is that relativistic quantum mechanics attempts to
associate localization with particles, whereas quantum field theory associates
localization with fields and observables.

This seemingly minor shift has profound consequences.

In relativistic quantum mechanics one seeks position operators,

\[
X^i,
\]

together with localized wavefunctions.

In quantum field theory one instead introduces operator-valued distributions,

\[
\phi(x),
\]

localized at spacetime points (or, more precisely, smeared over spacetime
regions), while particles appear only as particular excitations of these
fields.

Consequently, the ontology changes completely.

Particles become derived concepts.

Fields become fundamental.

The purpose of this chapter is to explain how this transition resolves many of
the conceptual difficulties encountered in relativistic quantum mechanics,
while simultaneously introducing new and equally profound questions regarding
vacuum structure, localization, and measurement.
